% 调用上级目录公共模板
\documentclass{../template/softeng}

\begin{document}

% 封面：文档名、项目名、版本、编写人、审核人、审批人、日期
\doccover{详细设计说明书}
{项目名称替换处}
{V2026.07}

% 目录
\tableofcontents
\newpage

% 修订记录
\begin{revrecord}
V1.0 & 2026-07-12 & 编写人 & 审核人 & 初稿完成，搭建完整详细设计国标文档框架 \\
\hline
\end{revrecord}

\section{引言}
\subsection{编写目的}
\subsection{项目背景}
\subsection{定义}
\subsection{参考资料}

\section{任务概述}
\subsection{需求概述}
\subsection{运行环境}
\subsection{条件与限制}

\section{模块详细设计}
\subsection{模块1设计}
\subsubsection{模块概述}
\subsubsection{模块输入输出接口}
\subsubsection{模块处理逻辑流程}
\subsubsection{模块内部数据结构}
\subsubsection{模块算法设计}
\subsubsection{模块异常处理逻辑}

\subsection{模块2设计}
\subsubsection{模块概述}
\subsubsection{模块输入输出接口}
\subsubsection{模块处理逻辑流程}
\subsubsection{模块内部数据结构}
\subsubsection{模块算法设计}
\subsubsection{模块异常处理逻辑}

\subsection{其他业务模块}

\section{接口详细设计}
\subsection{外部接口详细设计}
\subsubsection{用户交互接口}
\subsubsection{第三方软件接口}
\subsubsection{硬件接口}

\subsection{内部模块接口详细设计}
\subsubsection{模块间请求参数定义}
\subsubsection{返回数据结构定义}
\subsubsection{调用时序说明}

\section{数据结构详细设计}
\subsection{全局数据结构说明}
\subsection{数据库表详细设计}
\subsubsection{数据表1字段定义、约束、索引}
\subsubsection{数据表2字段定义、约束、索引}
\subsection{缓存/临时数据结构设计}

\section{算法详细设计}
\subsection{核心算法1}
\subsubsection{算法描述}
\subsubsection{输入输出参数}
\subsubsection{算法步骤流程图}
\subsubsection{伪代码/代码实现}
\subsubsection{边界条件处理}

\subsection{核心算法2}

\section{界面详细设计}
\subsection{页面1布局与控件说明}
\subsection{页面交互逻辑}
\subsection{页面输入校验规则}

\section{出错处理详细设计}
\subsection{模块内部异常分类}
\subsection{异常捕获流程}
\subsection{异常返回信息定义}
\subsection{异常恢复处理方案}

\section{安全详细设计}
\subsection{身份校验逻辑}
\subsection{数据加密处理设计}
\subsection{防攻击处理逻辑}

\section{维护详细设计}
\subsection{日志模块设计}
\subsection{配置项设计}
\subsection{数据备份逻辑}

% 流程图/界面截图插入示例
\begin{figure}[htbp]
    \centering
    % \includegraphics[width=0.7\textwidth]{module_flow.png}
    \caption{模块业务处理流程图}
\end{figure}

% 伪代码/代码块示例
\begin{lstlisting}[language=Python]
# 模块核心处理伪代码
def core_func(param):
    # 参数校验
    if check_param(param) is False:
        return error_result()
    # 业务逻辑处理
    data = query_db(param)
    res = calc_logic(data)
    return res
\end{lstlisting}

\end{document}